\section{Experimental Design and Evidence Selection}
\label{sec:setup}

We analyze the completed evaluations reported in the formal archive. RQ1 asks how integrated effectiveness varies by workload and exact goal overlap. RQ2 examines completed single-component removals and recorded binding behavior. RQ3 examines query-type and endpoint sensitivity in DBBench. RQ4 examines recorded inference and the cost of structural screening. All numerical results use completed records or saved timing measurements.

\subsection{Workloads, frozen memory, and overlap}
The primary workloads use 100 ScienceWorld \cite{wang2022scienceworld} test-split targets and indices 0--99 of the AgentBench DBBench workload \cite{liu2024agentbench}, executed through the project's local SQLite adapter. The archived study verifies the same target IDs for all eleven configurations within each workload. ScienceWorld allows 50 agent decisions and 100 engine steps. DBBench allows 30 decisions with a 15-round interface limit. These are the project protocols. All selected SQL targets are reads, so they do not exercise incorrect writes or transaction behavior.

The frozen memory contains 100 sources for each of ALFWorld \cite{shridhar2021alfworld}, ScienceWorld, and DBBench. Success/failure counts are 82/18, 53/47, and 79/21. Their source action lengths have medians 7, 24.5, and 2, respectively, and ranges 3--50, 4--50, and 2--7. The structural agent queries the 53 successful ScienceWorld sources or the 79 successful DBBench sources. Online updates are disabled.

Source and target task IDs are disjoint. After trimming outer whitespace and lowercasing, 36 ScienceWorld targets, covering 20 distinct goals, have a verbatim source goal match. We report the remaining 64 targets separately and retain the full set for the principal comparison. DBBench has no exact normalized goal overlap. Template overlap and semantic near duplicates may remain after these checks.

The archive also contains 100 ALFWorld training-slice targets at offset 3,253. We exclude that slice from the main comparisons. It is a development slice, and 99 of 100 archived structural-agent records include retrieved IDs from another benchmark. Fifty-two return 50 memories instead of five. The primary ScienceWorld and DBBench records consistently report five retrieved sources under the same-benchmark, success-only policy. The ALFWorld diagnostics therefore describe mixed versions of retrieval.

\subsection{Comparator scope}
The primary table includes No Memory, six memory adaptations, and integrated \SM{}. No Memory uses the base agent without retrieved guidance. Voyager-style \cite{wang2024voyager} uses lexical cosine over successful traces with compact skill guidance. Reflexion-style \cite{shinn2023reflexion} retrieves failed traces and disables online reflection updates. ExpeL-style \cite{zhao2024expel} may induce rules from retrieved successes and failures, with evaluation-time updates disabled. MemGPT-style \cite{packer2023memgpt} provides bounded archival recall. MCMA-style \cite{liang2026mcma} combines hybrid retrieval with diversified compact guidance. AgentKB-style \cite{tang2025agentkb} uses BM25 and a reliability filter with no updates. These styles are project adaptations. The class named \texttt{SemanticRetriever} in several of them uses token-count cosine.

A separate dense-only batch uses \texttt{all-MiniLM-L6-v2} to retrieve five successful, same-benchmark memories and runs the base retrieval agent. It has no \SM{} mapper, binder, or specialized execution routines. The complete secondary DBBench batch uses the same task IDs, frozen memory, and interface limits with endpoint label \texttt{grok-4.6}. We use that batch for the endpoint comparison. Incomplete secondary workloads and the later extension runs fall outside the selected evidence.

\subsection{Parameters and implementation version}
The inspected settings are $\alpha=0.55$ in Equation~\ref{eq:structural}, $(w_d,w_b,w_s)=(0.45,0.30,0.25)$ in Equation~\ref{eq:fusion}, and $\tau_p=0.75$ in Equation~\ref{eq:probe}. BM25 uses $k_1=1.5$ and $b=0.75$. These are fixed settings. The archived removals disable features or stages and cannot establish robustness across coefficients or thresholds.

The primary endpoint label is \texttt{gpt-5.6-sol} and the harness seed is 42. The client defaults to temperature zero and sends no API seed. Execution-time source commits and immutable model versions are unavailable. Sections~\ref{sec:design} and~\ref{sec:execution} describe the inspected code. The results apply to the logged variants. Agreement on the source eligibility policy and selected count supports the comparison at that boundary. Prompt, configuration, and executor identity at run time require their own version records.

\subsection{Terminal records and accounting}
The primary analysis uses one frozen \texttt{formal/episodes.jsonl} file. The formal study reports 3,300 terminal \texttt{completed} or \texttt{task-failed} records and 750 \texttt{agent-error} attempts across the three benchmark slices. The terminal records have one row per method and target. The additional attempts are associated with retries. We retain all valid terminal failures and use the selected archive for each comparison. The success rates describe terminal results under the archive's retry behavior. They do not measure first-attempt reliability.

Recorded tokens include context and action calls, plus usage restored from cached responses. JSON repair, retries, and SQLite caching affect the accounting. Reported totals measure logical model usage for terminal records. Billing, live latency, and retry overhead cannot be recovered from these summaries. Appendix~\ref{app:provenance} describes the material available to verify this version of the manuscript.
